\documentclass{article}
\usepackage[utf8]{inputenc}
\usepackage[spanish]{babel}
\usepackage{amsmath}
\usepackage{amssymb}
\usepackage{hyperref}
\usepackage{graphicx}
\usepackage[round, authoryear]{natbib}
\bibliographystyle{apalike} 
\usepackage{csquotes}
\usepackage[style=apa, backend=biber]{biblatex}
\addbibresource{references.bib}

\begin{document}



\subsection{Dinámica de la Vorticidad y la Enstrofía}
Para caracterizar cuantitativamente el desarrollo de la inestabilidad y la transición a la turbulencia en el régimen no lineal, es fundamental el estudio de la vorticidad y sus momentos estadísticos. En una configuración típica de KHI, el perfil inicial de velocidad transversal $V_y(x)$ se modela mediante una capa de cizalladura descrita por:

\begin{equation}
    V_y(x) \sim v_{\text{sh}} \tanh\left(\frac{x - x_0}{a_{\text{kh}}}\right),
    \label{eq:vel_profile}
\end{equation}

\noindent donde $v_{\text{sh}}$ representa la magnitud de la velocidad de cizalladura y $a_{\text{kh}}$ la escala característica del ancho de la capa.

\subsubsection{Cálculo Numérico y Contribución de Componentes}
En el contexto de la simulación numérica, la vorticidad se evalúa a partir de los gradientes discretos de la velocidad. La componente principal de la vorticidad en una geometría bidimensional es ortogonal al plano de simulación ($\omega_z$). La enstrofía asociada a esta componente, $\Omega_z(t)$, cuantifica la intensidad de los vórtices que rotan en el plano $xy$ y se define operativamente como:

\begin{equation}
    \Omega_z(t) = \int (\partial_x V_y - \partial_y V_x)^2 \, dA.
\end{equation}

Sin embargo, para capturar la dinámica completa del sistema, especialmente en configuraciones donde se permite la evolución de la tercera componente de la velocidad ($V_z$), es necesario considerar la enstrofía total $\Omega_{\text{tot}}$. Esta magnitud incorpora las contribuciones de los gradientes transversales de la velocidad fuera del plano (out-of-plane velocity), definidéndose como:

\begin{equation}
    \Omega_{\text{tot}}(t) = \int \left[ (\partial_y V_z)^2 + (\partial_x V_z)^2 + (\partial_x V_y - \partial_y V_x)^2 \right] \, dA.
\end{equation}

La comparación entre estas dos magnitudes permite evaluar la anisotropía de la turbulencia generada y el grado de acoplamiento entre los modos planares y la componente transversal. Para ello, se introduce el factor de fracción de enstrofía $f_{Vz}(t)$, que aísla la contribución relativa de los términos dependientes de $V_z$:

\begin{equation}
    f_{Vz}(t) = \frac{\Omega_{\text{tot}}(t) - \Omega_z(t)}{\Omega_{\text{tot}}(t)}.
\end{equation}

Un valor de $f_{Vz} \approx 0$ indica un flujo puramente bidimensional, mientras que un incremento en este factor señala la transferencia de energía hacia componentes ortogonales, un fenómeno relevante en plasmas magnetizados donde la tensión magnética puede inducir movimientos complejos fuera del plano de cizalladura.

\subsubsection{Enstrofía de Perturbación}
Dado que la enstrofía total está dominada inicialmente por la fuerte cizalladura del flujo base, es necesario aislar la contribución específica de las estructuras inestables emergentes. Siguiendo la metodología de \parencite{zrake-2016, takamoto-2011}, descomponemos la vorticidad instantánea $\omega_z$ restando el perfil promedio inicial. Definimos la vorticidad de perturbación $\omega'_z$ como:

\begin{equation}
    \omega'_z(x, y, t) = \omega_z(x, y, t) - \overline{\omega}_z^{(0)}(x),
\end{equation}

\noindent donde $\overline{\omega}_z^{(0)}(x) = \langle \omega_z(x, y, t=0) \rangle_y$ denota el promedio espacial a lo largo de la dirección longitudinal $y$ correspondiente al estado inicial. De esta forma, la enstrofía de perturbación se define como:

\begin{equation}
    \Omega'_z(t) = \int (\omega'_z)^2 \, dA.
\end{equation}

El análisis de $\Omega'_z(t)$ permite identificar las fases de crecimiento lineal y saturación no lineal, sirviendo como el principal diagnóstico para la tasa de crecimiento de la inestabilidad libre de la ``contaminación'' del flujo de fondo.

\printbibliography

\end{document}
